\chapter{Tooling and Reproducibility}
\label{app:tools}

This appendix documents the data-reduction and figure-generation tooling that
turns the raw Design Compiler reports into the tables and figures of
Chapter~\ref{ch:results}, so that the results are fully reproducible. The
complete source --- the modular accelerator, the DMR/TMR wrappers, the IJTAG
sensor-insertion tool, and the synthesis-sweep and data-reduction scripts ---
together with the thesis sources is available in the project repository at
\href{\repourl}{\repourl}.

\section{Directory layout}

The synthesis reports are organised by node and configuration:
\begin{itemize}
  \item \texttt{synth\_output/<node>/<config>/Report\_*.txt} --- the per-run
        area, cells, power, timing, references, and resources reports for the
        redundancy study (Scenarios A and B).
  \item \texttt{modules/out/<n>sensor/synthesis\_reports\_*/} --- the area, cells,
        power, and references reports for the IJTAG-instrumented builds, in
        \emph{counted} and \emph{black-box} accounting modes.
\end{itemize}

\section{Data-reduction scripts}

Three Python scripts perform the reduction; all are pure-Python plus
\texttt{matplotlib}/\texttt{numpy} and write into \texttt{thesis/data/} and
\texttt{thesis/img/plots/}:

\begin{itemize}
  \item \texttt{parse\_synth.py} --- parses every redundancy-study report,
        normalises power units to milliwatts, maps the on-disk configuration
        names to clean keys, and emits \texttt{synth\_data.csv} and a
        baseline-normalised \texttt{overhead\_scenarioB.csv}.
  \item \texttt{parse\_ijtag.py} --- parses the instrumented builds and emits
        \texttt{ijtag\_data.csv} and the counted-minus-black-box
        \texttt{ijtag\_cost.csv}.
  \item \texttt{make\_plots.py} and \texttt{make\_tables.py} --- generate every
        figure (as vector PDF) and every booktabs table directly from the CSVs.
\end{itemize}

Re-running the four scripts regenerates all numerical content of the thesis from
the raw reports. Listing~\ref{lst:power-parse} shows the core of the
unit-normalisation logic, which reads the power unit from each report header so
that the milliwatt/microwatt and picowatt/nanowatt differences between libraries
are reconciled automatically.

\begin{lstlisting}[language=mypython, caption={Power-unit normalisation in \texttt{parse\_synth.py} (excerpt).}, label={lst:power-parse}]
for ln in lines:
    m = re.search(r"Dynamic Power Units = 1(\w?)W", ln)
    if m:
        u = m.group(1)
        dyn_scale = {"m":1.0,"u":1e-3,"n":1e-6,"p":1e-9,"":1e3}.get(u, 1.0)
    m = re.search(r"Leakage Power Units = 1(\w?)W", ln)
    if m:
        u = m.group(1)
        leak_scale = {"m":1.0,"u":1e-3,"n":1e-6,"p":1e-9,"":1e3}.get(u, 1.0)
# top-of-hierarchy row holds the whole-design power (switch, int, leak, total)
\end{lstlisting}

\chapter{Configuration Reference}
\label{app:config}

\section{Redundancy enable constants (Scenario A)}

In the single-core flow the top-level design exposes the Boolean constants
\texttt{ADDERS\_DMR\_EN}, \texttt{ADDERS\_TMR\_EN}, \texttt{MULTIPLIERS\_DMR\_EN},
\texttt{MULTIPLIERS\_TMR\_EN}, \texttt{SHIFTERS\_DMR\_EN},
\texttt{SHIFTERS\_TMR\_EN}, \texttt{COMPARATORS\_DMR\_EN}, and
\texttt{COMPARATORS\_TMR\_EN}. Setting exactly one DMR/TMR constant true gives a
single-FU configuration; setting all four DMR (or all four TMR) constants true
gives full DMR (or full TMR); leaving all false gives the baseline.

\section{Replication constants (Scenario B)}

In the replicated flow, \texttt{replicate\_accl\_en} is enabled
(\texttt{ACCL\_NUM} = \texttt{FU\_NUM} = 3) and the redundancy lives in the
mapping layer, toggled by its \texttt{DMR\_ENABLED} and \texttt{TMR\_ENABLED}
constants: neither set gives the replicated baseline, \texttt{DMR\_ENABLED} gives
replicated-plus-DMR, and \texttt{TMR\_ENABLED} gives replicated-plus-TMR.

\section{IJTAG sensor configuration}

The IJTAG flow is driven by a single \texttt{config.json} that declares the
top-level entity, clock, libraries, packages, the activity-sensor type with its
ICL client-port names, and the sensor instances with their generic maps
(\texttt{SENSOR\_ID}, \texttt{WINDOW\_WIDTH}), functional port maps, and
hierarchical paths. The \texttt{configurations} section groups instances into the
cumulative \texttt{1sensor}/\texttt{2sensor}/\texttt{3sensor} builds used in
Section~\ref{sec:ij-synth}.
